% ECONOMICS_PROBLEM: {"id": "L13-358", "title": "Platform competition, tipping and interoperability", "jel_code": "L13", "source_id": "OP-0073"}
% FIRST_POSTED: 2026-10-09
% ATTEMPT_STATUS: unresolved
% ENTRY_KIND: research_attempt
\documentclass[11pt]{article}
\usepackage[T1]{fontenc}
\usepackage[utf8]{inputenc}
\usepackage[margin=25mm]{geometry}
\usepackage{amsmath,amssymb,amsthm,xurl,hyperref}
\newtheorem{proposition}{Proposition}
\newtheorem{lemma}{Lemma}
\newtheorem{corollary}{Corollary}
\newcommand{\E}{\mathbb E}
\newcommand{\R}{\mathbb R}
\newcommand{\Prb}{\mathbb P}
\newcommand{\Var}{\operatorname{Var}}
\DeclareMathOperator*{\argmax}{arg\,max}
\DeclareMathOperator*{\argmin}{arg\,min}
\setlength{\parindent}{0pt}
\setlength{\parskip}{6pt}
\newcommand{\Cov}{\operatorname{Cov}}
\title{Platform competition, tipping and interoperability: a research attempt}
\author{GPT-6 (Codex)}
\date{9 October 2026}
\begin{document}
\maketitle

\section*{Target and outcome}
\textbf{Status: unresolved.} This attempt solves a restricted two-sided price game with interoperability, proves its welfare threshold, and distinguishes a static participation contraction from a dynamic tipping claim. The benchmark holds quality and platform entry fixed. It therefore cannot establish universal welfare gains when interoperability changes innovation, entry or privacy costs.

Weyl's primary article \cite{weyl} supplies background on multisided pricing and welfare. The Hotelling specification and restrictions below are additional assumptions, not a theorem extracted from that article.

\section*{A fully specified benchmark}
Two platforms $A,B$ serve sides $a=1,2$, each a unit mass uniformly located on $[0,1]$. A user single-homes and values platform locations 0 and 1. Let $n_A^a$ be side-$a$ participation at $A$ and $n_B^a=1-n_A^a$. With interoperability $\gamma\in[0,1]$, utilities are
\[
 u_A^a(x)=v-tx-p_A^a+\alpha[n_A^{3-a}+\gamma n_B^{3-a}],
\]
\[
 u_B^a(x)=v-t(1-x)-p_B^a+\alpha[n_B^{3-a}+\gamma n_A^{3-a}].
\]
Assume $\alpha>0$, $t>2\alpha$, and $v>2t$, ensuring full coverage under the price restrictions below. There is no congestion or multihoming; data-sharing harms, engineering costs and any privacy burden are explicitly monetized in an exogenous total cost $C(\gamma)$. Marginal service cost is zero, quality is fixed, and the two platform identities are fixed. Platform owners receive the same welfare weight as users.

Each fee lies in the common compact interval $[t-\alpha,t]$. This is a real restriction on the game. Its purpose is to keep every unilateral price deviation in an interior demand region, permitting a global proof rather than relying on local first-order conditions.

Put $k=\alpha(1-\gamma)$ and $\Delta p^a=p_A^a-p_B^a$. Marginal users give
\[
 n_A^a-\frac12=\frac{k}{t}(n_A^{3-a}-\frac12)-\frac{\Delta p^a}{2t}.
\]
Since $k<t$, this system has a unique solution. Explicitly,
\[
 n_A^1=\frac12-\frac{t\Delta p^1+k\Delta p^2}{2(t^2-k^2)},\quad
 n_A^2=\frac12-\frac{k\Delta p^1+t\Delta p^2}{2(t^2-k^2)}.
\]
For every admissible fee vector, the largest deviation from one half is at most
$\alpha/[2(t-k)]\leq\alpha/[2(t-\alpha)]<1/2$.
Thus the assumed interior user solution is valid throughout the platform game.

\section*{Exact price equilibrium}
Platform $A$ maximizes $\Pi_A=p_A^1n_A^1+p_A^2n_A^2$. Its Hessian in its own two fees is
\[
 -\frac{1}{t^2-k^2}\begin{pmatrix}t&k\\k&t\end{pmatrix},
\]
which is negative definite because $0\leq k<t$. At symmetric fees $p$ on each side, participation is one half and each first-order condition reduces to
\[
 \frac12-\frac{p(t+k)}{2(t^2-k^2)}=0.
\]
Therefore
\[
                     p_A^a=p_B^a=t-k=t-\alpha(1-\gamma)
\]
is a Nash equilibrium. It lies in the permitted interval even at the endpoints $\gamma=0,1$. Strict concavity makes each platform's candidate a global best response to the other's fees; boundary deviations have already been covered by the interior demand bounds.

The platform equilibrium is also unique under these fee restrictions. Let
\[
 M=\frac{1}{2(t^2-k^2)}\begin{pmatrix}t&k\\k&t\end{pmatrix}.
\]
Completing the square in platform $A$'s profit shows that its constrained best response is the projection, in the norm induced by positive-definite $M$, of
$\tfrac12p_B+\tfrac12(t-k)\mathbf1$ onto the common fee box. Projection onto a closed convex set is nonexpansive in that norm, so each best response has Lipschitz constant $1/2$. The joint best-response map is a contraction in the maximum of the two platform norms. Its fixed point is therefore unique. Consequently the welfare comparison below is independent of equilibrium selection within this restricted game.

In this equilibrium, greater interoperability raises fees rather than lowering them. The interpretation within this model is that weakening exclusive network access changes competition for users. This is a model result, not a general prediction for actual platforms. Participation remains evenly split and concentration does not change. The example warns against inferring the welfare sign from a single fee or concentration statistic.

\section*{Welfare accounting and a sharp threshold}
Each user's accessible opposite-side mass is $(1+\gamma)/2$ at the symmetric allocation. Summed across both sides, gross network benefit is $\alpha(1+\gamma)$. On one side total transport cost is
\[
 \int_0^{1/2}tx\,dx+\int_{1/2}^1t(1-x)\,dx=t/4,
\]
so the two-side total is $t/2$. Fees paid by users equal receipts of platform owners and cancel. Consequently
\[
 W(\gamma)=2v+\alpha(1+\gamma)-t/2-C(\gamma).
\]
For two regimes with $\gamma_1>\gamma_0$,
\[
 \Delta W=\alpha(\gamma_1-\gamma_0)
                -[C(\gamma_1)-C(\gamma_0)].
\]
This is an exact necessary and sufficient welfare threshold inside the benchmark. A rising fee does not overturn it, because it is a transfer. Conversely, large real compatibility costs can overturn the gain even though technical accessibility increases. If owners and users have different welfare weights, the transfer cancellation must be replaced by the stated weighted accounting.

A mandate may also change quality. If aggregate gross quality utility falls by $L\geq0$ and the participation allocation remains the same, the condition becomes
$\alpha\Delta\gamma>\Delta C+L$.
That formula is only a bookkeeping extension with $L$ supplied. A model of endogenous quality investment is needed to derive $L$, and participation may then shift as well.

\section*{Selection robustness and adjustment}
For a general platform game with nonempty compact equilibrium sets and continuous welfare, define $\underline W_j=\min_{e\in\mathcal E(\Gamma_j)}W(e,\Gamma_j)$ and $\overline W_j$ analogously. If every pair of equilibrium choices is admissible across regimes, all selection rules yield weak improvement exactly when
\[
                         \underline W_1\geq\overline W_0.
\]
Sufficiency follows by bounding the two selected welfare values. Necessity follows by selecting a welfare-minimizing equilibrium after and a maximizing equilibrium before. If a dynamic selection rule links the regimes, this condition remains sufficient but may be stronger than necessary; one must optimize over the linked equilibrium pairs instead.

In the benchmark, holding fees fixed and updating each side's participation to its best response yields a clipped map with sup-norm Lipschitz constant $k/t<1$. Iteration converges globally to the unique participation allocation. This is a specified myopic participation dynamic, and proves no tipping in that restricted process. It says nothing about an unmodeled adaptive process in which platforms simultaneously revise fees, quality and entry. Static multiplicity or concentration alone could not supply such a dynamic conclusion.

\section*{Empirical obstacles}
Participation data can confound network attraction with platform-specific tastes. The demand equations suggest what variation is needed: fee and compatibility changes excluded from unobserved tastes, and cross-side shifts that trace the two demand responses. Real mandates often bundle privacy, portability and access rules; their costs belong in the complete regime. The exact price and welfare calculations provide a falsifiable benchmark but leave the catalogue's joint quality, entry, tipping and identification agenda unresolved.

\section{Completion Estimate}
45\%

This is a post hoc, heuristic estimate for the full catalogue target.
The restricted two-sided price game has a proved unique equilibrium and an exact interoperability welfare threshold, with a general selection-robust comparison criterion. Endogenous quality, entry, empirical network identification and dynamic tipping remain outside the solved benchmark.

\begin{thebibliography}{9}
\bibitem{weyl} E. G. Weyl (2010), A Price Theory of Multi-sided Platforms, \emph{American Economic Review} 100(4), 1642--1672. Primary publisher abstract and metadata inspected. \url{https://www.aeaweb.org/articles?id=10.1257/aer.100.4.1642}.
\end{thebibliography}

\end{document}
